\section{Conclusion and Future Work}
\label{sec:conclusion}

We introduced DyFlow-JEPA, a self-supervised framework that reformulates unsupervised anomaly detection in multi-sensor IoT and remote sensing telemetry as vector field modeling on a regularized latent dynamical manifold. By avoiding raw-space ODE integration and eliminating signal-level reconstruction traps, DyFlow-JEPA models physical systems as continuous tangent velocity fields. Non-contrastive manifold regularization enforces representation diversity while decorrelating latent dynamical modes across coupled sensor channels. Coupled with a single-pass Deterministic Midpoint Velocity discrepancy metric ($t=0.5$), Mahalanobis covariance whitening, and Extreme Value Theory (EVT / SPOT) tail calibration, DyFlow-JEPA computes calibrated, threshold-free anomaly alarms with low inference latency ($3.47$s).

Evaluations across 11 benchmark telemetry datasets comprising 130 continuous channels spanning four IoT and remote sensing domains (and extended distributed streaming ablations on Exathlon, totaling over 1,300 model runs) demonstrate that DyFlow-JEPA outperforms baseline contrastive and reconstructive architectures (TimesNet, TranAD, Anomaly Transformer, DCdetector, and NCAD) across both channel-weighted and unweighted dataset macro metrics. Ablation results confirm that Optimal Transport Flow Matching mitigates conditional-mean collapse on multi-branching dynamics and provides stability against non-stationary streaming drift.

From an operational perspective on remote satellite downlinks and edge IoT monitoring, while strict timestep-level Point-F1 scores in the $0.20-0.35$ range reflect the low base-rate and stochastic noise of continuous sensor streams, DyFlow-JEPA's $+52.5\%$ relative improvement over reconstructive baselines reduces false alarms and provides a more reliable prioritization signal for operators.

Future research directions include:
\begin{itemize}[noitemsep, topsep=0pt, leftmargin=*]
    \item \textit{Differentiable Physical Constraint Layers:} For physical domains where governing laws are known \textit{a priori} (e.g., Kirchhoff circuit laws in power grids or kinematics in robotics), integrating differentiable physical constraint layers directly into the DyFlow-JEPA latent velocity field to ensure boundary compliance.
    \item \textit{Continuous Online Streaming Adaptation:} Developing continuous test-time momentum updates to track gradual seasonal or orbital regime drifts without catastrophic forgetting.
    \item \textit{Edge Hardware Quantization:} Deploying lightweight integer-quantized (INT8) DyFlow-JEPA encoders directly on radiation-hardened spacecraft microcontrollers and wearable biomedical sensors.
    \item \textit{Foundation Models for Cyber-Physical Telemetry:} Pretraining multi-domain DyFlow-JEPA world models on multi-mission telemetry archives for zero-shot downstream fault diagnosis.
\end{itemize}

